\documentclass[10.5pt]{article}
\usepackage[letterpaper, margin=0.6in, top=0.5in, bottom=0.5in]{geometry}
\usepackage{amsmath,amssymb,amsthm}
\usepackage{fancyhdr}
\usepackage{xcolor}
\usepackage{tikz}
\usepackage{enumitem}
\usepackage{tcolorbox}
\usepackage{hyperref}

% Define colors matching instructor board annotations
\definecolor{chalkblue}{RGB}{0, 105, 185}
\definecolor{chalkred}{RGB}{205, 35, 35}
\definecolor{chalkpink}{RGB}{230, 90, 120}
\definecolor{chalkcyan}{RGB}{15, 160, 190}
\definecolor{chalkgreen}{RGB}{20, 130, 50}
\definecolor{boxbg}{RGB}{252, 252, 254}
\definecolor{boxframe}{RGB}{150, 160, 180}

\pagestyle{fancy}
\fancyhf{}
\renewcommand{\headrulewidth}{0pt}
\renewcommand{\footrulewidth}{0pt}

\setlength{\parindent}{0pt}
\setlength{\parskip}{2.5pt}

\begin{document}

% ==============================================================================
% PAGE 1
% ==============================================================================
\noindent \textbf{\large Math 105L} \hfill \textbf{\large Lecture 5-3}\\[1pt]
{\textbf{\Large The Second Derivative, Concavity, \& Inflection Points}}\\[-5pt]
\rule{\textwidth}{1.2pt}

\vspace{0.03in}

\textbf{\large 1. Review: What the First Derivative Tells Us}

The first derivative $f'(x)$ gives the instantaneous rate of change of $f$ (slope of the tangent line):
\begin{itemize}[leftmargin=1.2em, topsep=1pt, itemsep=2pt]
  \item If $\mathbf{f'(x) > 0}$ on an interval, then $f(x)$ is \textbf{increasing} on that interval.
  \item If $\mathbf{f'(x) < 0}$ on an interval, then $f(x)$ is \textbf{decreasing} on that interval.
\end{itemize}

\vspace{0.03in}

\textbf{\large 2. The Second Derivative: Rate of Change of Slope}

The \textbf{second derivative} is simply the derivative of the first derivative:
\[
  \mathbf{f''(x) = \frac{d}{dx}\big[f'(x)\big] = \frac{d^2 f}{dx^2} = \lim_{h\to 0}\frac{f'(x+h)-f'(x)}{h}}
\]
Just as $f'(x)$ tells us whether $f(x)$ is increasing or decreasing, $f''(x)$ tells us whether the \textbf{slopes $f'(x)$ themselves} are increasing or decreasing.

\vspace{0.03in}

\begin{tcolorbox}[colback=boxbg, colframe=boxframe, boxrule=0.8pt, arc=3pt, left=8pt, right=8pt, top=4pt, bottom=4pt]
\textbf{\large Concavity: Geometric \& Analytic Characterizations}
\begin{itemize}[leftmargin=1.2em, topsep=2pt, itemsep=5pt]
  \item \textbf{Concave Up ($\cup$):}
  \begin{itemize}[label=--]
    \item \textbf{Slope Behavior:} Slopes are \textbf{increasing} as $x$ moves left to right ($f''(x) > 0$).
    \item \textbf{Tangent Line Property:} The tangent lines lie \textbf{strictly below} the curve.
    \item \textbf{Intuition:} The graph holds water like a cup.
  \end{itemize}
  \item \textbf{Concave Down ($\cap$):}
  \begin{itemize}[label=--]
    \item \textbf{Slope Behavior:} Slopes are \textbf{decreasing} as $x$ moves left to right ($f''(x) < 0$).
    \item \textbf{Tangent Line Property:} The tangent lines lie \textbf{strictly above} the curve.
    \item \textbf{Intuition:} The graph sheds water like an umbrella.
  \end{itemize}
\end{itemize}
\end{tcolorbox}

\vspace{0.04in}

\textbf{\large 3. Crucial Subtlety \& Counterexample: $f(x) = x^4$}

A common misconception is that concavity requires $f''(x) > 0$ strictly everywhere or that $f''(c) = 0$ implies a change in concavity. Consider:
\[
  f(x) = x^4 \implies f'(x) = 4x^3 \implies \mathbf{f''(x) = 12x^2}
\]
\begin{itemize}[leftmargin=1.2em, topsep=1pt, itemsep=3pt]
  \item At $x = 0$, we have $f''(0) = 12(0)^2 = 0$.
  \item For all $x \ne 0$, $f''(x) = 12x^2 > 0$.
  \item The slopes $f'(x) = 4x^3$ are strictly increasing across all $x \in \mathbb{R}$ (from negative, through $0$, to positive).
  \item The horizontal tangent line at $(0,0)$ ($y=0$) lies strictly below the curve $y=x^4$ for all $x \ne 0$.
  \item \textbf{Conclusion:} $f(x) = x^4$ is \textbf{strictly concave up everywhere on $(-\infty, \infty)$}, even though $f''(0) = 0$.
\end{itemize}

\vspace{0.04in}

\textbf{\large 4. Inflection Points}

\begin{tcolorbox}[colback=chalkblue!6, colframe=chalkblue, boxrule=1pt, arc=3pt, left=8pt, right=8pt, top=4pt, bottom=4pt]
\textbf{Definition of an Inflection Point:}\\[2pt]
A point $(c, f(c))$ on the curve is an \textbf{inflection point} if:
\begin{enumerate}[leftmargin=1.4em, topsep=1pt, itemsep=2pt]
  \item The function $f$ is continuous at $c$ and has a tangent line at $x=c$.
  \item The concavity of $f$ \textbf{changes sign} across $x=c$ (from concave up to concave down, or vice versa).
\end{enumerate}
\textbf{Geometric Signature:} At an inflection point, the \textbf{tangent line cuts across (crosses through)} the curve!
\end{tcolorbox}

\vspace{0.03in}

\textbf{Test for Inflection Points:}
\begin{enumerate}[leftmargin=1.2em, topsep=1pt, itemsep=2pt]
  \item Find candidate points where $f''(x) = 0$ or $f''(x)$ does not exist.
  \item Test test values on intervals to see if $f''(x)$ actually \textbf{changes sign} (e.g. from $+$ to $-$ or $-$ to $+$).
  \item If $f''(x)$ does not change sign (as in $f(x)=x^4$ where $f''$ is positive on both sides of $0$), it is \textbf{not} an inflection point!
\end{enumerate}

\end{document}

